\SourceLeaf{48}
als auch der zweite Teil des vorhergehenden Gliedes $(r-1)^{\text{ti}}$,
\[
= \frac{1}{2^{2m-1}}\,{}^{2m-1}\mathfrak{B}^{(m+r-1)}\cos 2r\varphi,
\]
mit $\cos 2r\varphi$ multipliziert sind. Die übrigen Glieder ergeben kleinere oder größere Vielfache. Fasst man also diese Glieder zusammen, so erhält man in der neuen Reihe für $\cos\varphi^{2m}$ das allgemeine Glied $r^{\text{tus}}$. Wegen
\[
{}^{2m-1}\mathfrak{B}^{(m+r)} + {}^{2m-1}\mathfrak{B}^{(m+r-1)} = {}^{2m}\mathfrak{B}^{(m+r)},
\]
nimmt es die Form an
\[
\frac{1}{2^{2m-1}}\,{}^{2m}\mathfrak{B}^{(m+r)} . \cos 2r\varphi.
\]

\par\noindent Dieser Ausdruck stimmt mit dem unten für die geraden Potenzen angenommenen Gesetz vollständig überein. Daraus folgt, dass dieses Gesetz für eine gerade Potenz gilt, sobald es für die nächstniedrigere ungerade Potenz bewiesen ist.

4. Nun multipliziere man die eben für $\cos\varphi^{2m}$ hergeleitete Reihe
\[
\cos\varphi^{2m} = \frac{1}{2^{2m-1}}\frac{{}^{2m}\mathfrak{B}^{(m)}}{2}
+ \frac{1}{2^{2m-1}}\,{}^{2m}\mathfrak{B}^{(m+1)}\cos 2\varphi
+ \frac{1}{2^{2m-1}}\,{}^{2m}\mathfrak{B}^{(m+2)}\cos 4\varphi + .\,.\,.
\]
\[
{} + \frac{1}{2^{2m-1}}\,{}^{2m}\mathfrak{B}^{(m+r)}\cos 2r\varphi + .\,.\,.
\]
mit $\cos\varphi$. Auf der einen Seite erhält man so $\cos\varphi^{2m+1}$, auf der anderen in jedem Glied außer dem Anfangsglied ein Produkt zweier Kosinus, das wie zuvor in zwei Teile zu zerlegen ist. So liefert das erste Glied
\[
\frac{1}{2^{2m-1}}\,{}^{2m}\mathfrak{B}^{(m+1)}\cos 2\varphi\cos\varphi
\]
die Teile
\[
\frac{1}{2^{2m-1}}\,{}^{2m}\mathfrak{B}^{(m+1)}\frac{1}{2}\cos\varphi
+ \frac{1}{2^{2m-1}}\,{}^{2m}\mathfrak{B}^{(m+1)}\frac{1}{2}\cos 3\varphi.
\]

\par\noindent Addiert man seinen ersten Teil zu dem einzigen Teil, den das Anfangsglied liefert,
\[
\frac{1}{2^{2m-1}}\frac{{}^{2m}\mathfrak{B}^{(m)}}{2}\cos\varphi,
\]
so erhält man das Anfangsglied der neuen Reihe
\[
\frac{1}{2^{2m}}\,{}^{2m+1}\mathfrak{B}^{(m+1)}\cos\varphi.
\]
